% !TeX program = xelatex
% Compile twice with XeLaTeX. The numbered displays and census entries are shared by both editions.
\documentclass[UTF8,a4paper,11pt,scheme=plain]{ctexart}
\usepackage[left=26mm,right=26mm,top=23mm,bottom=24mm,footskip=11mm]{geometry}
\usepackage{amsmath,amssymb,amsthm}
\usepackage[unicode,hidelinks]{hyperref}
\hypersetup{pdftitle={三元域前缀主子式最短路的仿射像约化 / An affine-image quotient for ternary prefix minors},pdfauthor={}}
\ctexset{section={format=\large\bfseries\raggedright,name={},number=\arabic{section},
  aftername=\quad,beforeskip=2.4ex plus .3ex,afterskip=.9ex plus .2ex}}
\newtheoremstyle{zh}{.55\baselineskip}{.38\baselineskip}{\normalfont}{}{\bfseries}{}{.7em}{}
\theoremstyle{zh}
\newtheorem{zhthm}{定理}
\theoremstyle{plain}
\newtheorem{enthm}{Theorem}
\linespread{1.07}
\setlength{\emergencystretch}{1.2em}
\setlength{\parindent}{2em}
\newcommand{\F}{\mathbb F_3}
\newcommand{\A}{\mathcal A}
\newcommand{\im}{\operatorname{im}}
\newcommand{\GL}{\operatorname{GL}}
\newcommand{\stateeq}{u^\top b=1,\qquad u^\top W=0,\qquad \A(s)=b+\im W.}
\newcommand{\minoreq}{\det (I-u_B z_B^\top)=1-\sum_{i\in B}u_i z_i.}
\newcommand{\actioneq}{(u,W,b)\longmapsto(Pu,P^{-\top}W,P^{-\top}b),\qquad
       (u,\A)\longmapsto(Pu,P^{-\top}\A).}
\newcommand{\frameeq}{(W,b)\longmapsto(WT,b+Wt),\qquad T\in\GL_d(\F),\quad t\in\F^d.}
\newcommand{\homoeq}{H(W,b)=\im\begin{pmatrix}W&b\\0&1\end{pmatrix}\subseteq \F^{m+1},
       \qquad (z,1)\in H(W,b)\ \Longleftrightarrow\ z\in\A(s).}
\newcommand{\fivewitness}{R_0\gets R_0+R_1,\qquad R_0\gets R_0+R_2,\qquad R_3\gets R_3-R_4.}
\newcommand{\counttable}[6]{%
  \begin{center}
  \begingroup\small\setlength{\tabcolsep}{8pt}\renewcommand{\arraystretch}{1.18}
  \begin{tabular}{@{}lrrr@{}}
  \hline
  #1 & #2 & #3 & #4\\ \hline
  #5 & 58 & 26 & 24\\
  #6 0 & 1 & 1 & 1\\
  #6 1 & 6 & 6 & 6\\
  #6 2 & 40 & 40 & 35\\
  #6 3 & 166 & 126 & 80\\ \hline
  \end{tabular}
  \endgroup
  \end{center}}
\newcommand{\sources}[1]{%
  \begin{thebibliography}{9}\small
  \bibitem{#1orbits}
  SHIN09-git. Monomial guard settlement, \texttt{ORBITS.md}.
  \href{https://github.com/SHIN09-git/integer-mult-bounds/blob/03caa7f6cb5bd21c834024e9d81f6b43eb325b91/research/monomial-guard-settlement/ORBITS.md}{Source at \texttt{03caa7f}.}
  \bibitem{#1lean}
  SHIN09-git. Exact shortest paid paths for normalized full affine guards over $\F$.
  \href{https://github.com/SHIN09-git/integer-mult-bounds/blob/5703e56d6eb4b4f0a83af8c23a5d1dddd0b5474f/research/monomial-guard-orbit-lean/README.md}{Lean package at \texttt{5703e56}.}
  \bibitem{#1charge}
  N.~Petalas. Charge-only ternary guard costs through two additions.
  \href{https://github.com/CrocSwap/integer-mult-bounds/pull/6793}{Project contribution 6793.}
  \end{thebibliography}}

\begin{document}
\fontsize{12}{16.5}\selectfont
\renewcommand{\proofname}{证明}
\renewcommand{\refname}{参考文献}
\pdfbookmark[0]{中文}{edition-zh}
\begin{center}
  {\sffamily\bfseries\fontsize{20}{26}\selectfont 三元域前缀主子式最短路的仿射像约化\par}
  \vspace{.5em}
  {\small 研究短稿\quad 2026 年 11 月 21 日\qquad\hyperlink{edition-en.0}{English version}}
\end{center}
\noindent\textbf{摘要}\quad
在三元域上，对一族秩一扰动矩阵，目标是用最少的收费行加法，使同一行序下所有参数的真前缀主子式都非零。
现有算法已合并免费行置换和逐行变号。我们证明，还可以用对偶向量的仿射像代替它的参数表示。
给定原始列，相同仿射像保留每个操作序列的成功判定，从商图仍能恢复在原行号上可执行的最优指令。
在原五行算例中，距起点恰好三步的状态类从 166 个减至 80 个，最优代价仍为 3。

\section{模型与最短路定理}
记 $K=\F$，行号为 $0,\ldots,m-1$，其中 $m\ge1$。一个带行号的状态为
$s=(u,W,b)$，其中 $u,b\in K^m$、$W\in K^{m\times d}$。
参数 $a$ 遍历整个 $K^d$，对应的对偶列为 $z=b+Wa$。原问题的归一化条件与这些列组成的仿射像分别写为
\begin{equation}\label{z:state}\stateeq\end{equation}
对行指标集 $B$，秩一行列式恒等式给出
\begin{equation}\label{z:minor}\minoreq\end{equation}
给定一个行序 $\pi$，取其前 $k$ 项的指标集 $B_k^\pi$，其中 $0\le k<m$。
若对每个 $z\in\A(s)$ 和每个这样的 $k$，式~\eqref{z:minor} 均非零，则称 $\pi$ 通过检验。
空前缀的行列式为 $1$。整矩阵的行列式由归一化条件恒为零，故只检验真前缀。
成功要求存在一个对全族同时有效的行序。

行置换与逐行乘以非零标量免费，在 $K$ 上后者就是逐行变号。
每次 $R_i\gets R_i+cR_j$ 收费 $1$，其中 $i\ne j$，$c\in\{1,-1\}$。
令 $P$ 为该操作在原始列 $u$ 上的矩阵，记 $P^{-\top}=(P^{-1})^\top$。
状态和仿射像按下式变换。
\begin{equation}\label{z:action}\actioneq\end{equation}
例如，收费加法更新 $u_i\gets u_i+cu_j$，并对每个对偶列更新 $z_j\gets z_j-cz_i$。
记 $C(s)$ 为从 $s$ 到任一成功状态所需的最少收费次数，免费操作可以穿插在任意位置。

\begin{zhthm}\label{z:main}
设两个带行号的状态有相同的 $u$ 和 $\A$，其参数维数可以不同。对任意合法操作序列及任意指定的最终行序，
从两个状态执行该序列得到相同的成功判定。因此它们到成功集合的最短距离相等。
进一步按免费行置换和变号合并这些 $(u,\A)$ 后，商图仍给出 $C(s)$，并可恢复原输入上的最优指令。
\end{zhthm}
\noindent\emph{证明。}\enspace
同一操作按式~\eqref{z:action} 将相等的两对 $(u,\A)$ 送到相等的两对。
对序列长度归纳，再用式~\eqref{z:minor} 检查给定行序的每个前缀，即得第一句。
两边能成功的指令及其收费次数相同。免费行操作还可能改变行号，需要另行处理。

\newpage
\section{免费操作、参数换元与指令恢复}
令 $Q$ 为带符号的置换矩阵。若两个代表的 $(u,\A)$ 相差 $Q$ 的作用，
则在第二个代表上做加法 $E=I+ce_i e_j^\top$，等价于先在第一个代表上做 $Q^{-1}EQ$，
再施加同一个 $Q$。共轭后的矩阵仍是一次合法行加法，只有行号和系数可能改变。
免费操作自身留在同一类内。若第二个代表有通过检验的行序，把行号换回便得到第一个代表的可用行序。
这一论证只用到仿射像的作用，因而也适用于参数写法不同的代表。

在商图中以收费加法作单位长度的有向边，允许自环。原执行投影为同收费长度的游走。反过来，
上一段的共轭关系允许从任一实际代表逐边执行，并在终点选择一个通过检验的免费行序。
因此商图到成功集合的距离就是 $C(s)$。按层搜索时，每个新类保存一个从原输入到达的带行号状态、
它的父类和最后一次加法。沿父指针反向读取并倒序排列，即得实际输入上的最短指令。
这完成了定理~\ref{z:main} 的证明。\qed

\medskip\noindent\textbf{仿射像的判定。}
公共换元 $a\mapsto Ta+t$ 在数据上给出
\begin{equation}\label{z:frame}\frameeq\end{equation}
它保持仿射像。反之，在参数维数 $d$ 相同时，如果 $b+\im W=b'+\im W'$，
就有 $\im W=\im W'$ 且 $b'-b=Wt$。把 $\ker W'$ 同构到 $\ker W$，
再把两者的补空间通过共同的像对应起来，所得可逆映射 $T$ 满足 $W'=WT$。
该论证允许斜率列线性相关，也包括 $d=0$。

不同参数维数也可直接比较。引入齐次坐标，令
\begin{equation}\label{z:hom}\homoeq\end{equation}
最后坐标为 $1$ 的切片张成整个 $H$。事实上，两个切片向量的差遍历
$\im W\times\{0\}$，再添上 $(b,1)$ 即可。因此，固定行号时，$u$ 与上述分块矩阵转置的
简化行阶梯形中的非零行，完整刻画了 $(u,\A)$。如需再合并免费操作，可对每个 $Q$
计算 $Qu$ 与 $\operatorname{diag}(Q^{-\top},1)H$ 的编码，再取字典序最小值。

对一个参数，式~\eqref{z:frame} 只有六种换元 $a\mapsto\pm a+t$。
对每种换元，把每个完整行三元组 $(u_i,w_i,b_i)$ 与它的负值取较小者，再排序全部行并保留重数。
六种结果中的最小者就是本文复算使用的编码。斜率和常数项在 $u_i=0$ 的行上也保留。

\medskip\noindent\textbf{指定终点时的区别。}
取 $u=(1,0)^\top$、$W=(0,1)^\top$，分别令 $b=(1,0)^\top$ 和 $b'=(1,1)^\top$。
两者的仿射像都是 $\{(1,a)^\top:a\in K\}$，因而在新商图中是同一顶点。
但非零的 $u$ 固定第一行，逐行变号也无法把第二行的常数从 $0$ 变为 $1$。
所以两个指定的原始状态之间没有零成本路径。

\newpage
\section{精确枚举与适用范围}
旧算法用完整行的多重集表示免费等价类，并已在 Lean 中证明其最短路和指令恢复正确\cite{zorbits,zlean}。
另有一个只用电荷的结果判定至多两步的代价\cite{zcharge}。本稿保留原始列和整个对偶仿射像，
所证的最短路结论适用于任意步数。

下表比较三种编码。每列都先合并免费行置换和变号，后两列还分别合并参数平移及全部可逆仿射换元。
五行例的第 $k$ 层按各商图中到起点的\emph{最短}距离恰为 $k$ 计数。
\counttable{计数范围}{仅免费行操作}{加参数平移}{加全部换元}{三行全部归一化输入}{五行例，距离 }

三行、一个参数时，归一化输入共 $26\cdot9\cdot9=2\,106$ 个。
有 $26$ 个非零的 $u$ 可选，内积条件分别为 $W$ 和 $b$ 留下 $9$ 个选择。
脚本对原始带行号的图作多源 $0$--$1$ 广度优先搜索，显式枚举免费交换、变号和收费加法，
并直接计算 $I-uz^\top$ 的各前缀行列式。商图程序对每个前缀 $B$ 检查
$u_B^\top W_B=0$ 和 $u_B^\top b_B\ne1$，因为非零线性泛函取遍 $K$。将两种方法所得距离逐个比较，
全部吻合。最优代价为 $0$ 的有 $882$ 个，为 $1$ 的有 $1\,080$ 个，为 $2$ 的有 $144$ 个。

五行例沿用前序包的 $u=(1,1,1,1,1)^\top$、$W=(1,1,1,1,-1)^\top$、
$b=(1,1,1,1,0)^\top$。搜索在第 3 层首次成功，恢复出的零起始行号指令为
\[
  \fivewitness
\]
执行后，对每个参数都有 $(u_i z_i)_i=(0,0,0,0,1)$，故原行序的每个真前缀主子式均为 $1$。
第三层少了 $86$ 个状态类。表中给出的是层大小，实际运行时间还取决于编码成本。

另对两行、$d=0,1,2$ 的全部 $819$ 个 $(W,b)$，脚本比较式~\eqref{z:hom} 的编码和逐点生成的仿射像。
两者给出完全相同的 $22$ 类。这项检查针对像集编码，不加归一化条件。
复算使用 Python 3.12.10 标准库，在本研究包目录执行 \texttt{python -B verify.py} 即可。

\medskip\noindent\textbf{形式化与后续问题。}
Lean 4.24.0 已验证定理~\ref{z:main} 的任意长度路径提升、最优指令恢复和五行例的三步最优性。
Lean 还验证了五行例的闭球计数。半径从二增至三时，按免费行操作归类的类数由 $47$ 增至 $213$，
进一步按仿射像归类的类数由 $42$ 增至 $122$。取差得到第三层的 $166$ 和 $80$ 类。表中各数均可用 Python 复算。
多参数时，两种编码的耗时仍待比较。
\sources{z}

\clearpage
\normalsize
\pdfbookmark[0]{English}{edition-en}
\setcounter{section}{0}\setcounter{equation}{0}
\renewcommand*{\theHsection}{en.\arabic{section}}
\renewcommand*{\theHequation}{en.\arabic{equation}}
\renewcommand{\proofname}{Proof}
\renewcommand{\refname}{References}
\setlength{\parindent}{1.5em}
\begin{center}
  {\sffamily\bfseries\fontsize{17}{22}\selectfont
  An affine-image quotient for shortest paths\\
  to nonvanishing prefix minors over $\F$\par}
  \vspace{.5em}
  {\small Research note\quad 21 November 2026\qquad\hyperlink{edition-zh.0}{Chinese version}}
\end{center}
\noindent\textbf{Abstract.}\quad
For a family of rank-one perturbations of the identity over $\F$, we seek the fewest paid row
additions that make every proper leading principal minor nonzero in one order common to all
parameters. The existing search quotients out free row permutations and sign changes. We show
that the parametrized dual vector can also be replaced by its affine image. With the primal
column fixed, equal images give the same success test after every operation word. An optimal
word can be recovered on the original row labels. In a five-row example the number of classes
at distance exactly three from the start drops from 166 to 80, and the optimum remains three.

\section{The model and the shortest-path theorem}
Write $K=\F$ and index the $m\ge1$ rows from $0$ to $m-1$. A labeled state is
$s=(u,W,b)$ with $u,b\in K^m$ and $W\in K^{m\times d}$. The dual column is $z=b+Wa$
as $a$ runs over all of $K^d$. We assume the usual normalization and write $\A(s)$ for the
set of dual columns, so that
\begin{equation}\label{e:state}\stateeq\end{equation}
For an index set $B$, the rank-one determinant identity gives
\begin{equation}\label{e:minor}\minoreq\end{equation}
For an order $\pi$, let $B_k^\pi$ consist of its first $k$ indices, $0\le k<m$.
The order passes if~\eqref{e:minor} is nonzero for every $z\in\A(s)$ and every such $k$.
The empty minor is $1$. The determinant of the full matrix is zero by normalization.
Success means that one order passes for the entire family.

Row permutations and independent row scalings by nonzero field elements are free. Over $K$
the scalings are sign changes. Each addition $R_i\gets R_i+cR_j$ with $i\ne j$ and
$c\in\{1,-1\}$ costs one. Let $P$ denote the operation matrix on the primal column $u$ and
put $P^{-\top}=(P^{-1})^\top$. The actions on the presentation and on the affine image are
\begin{equation}\label{e:action}\actioneq\end{equation}
For an addition, $u_i$ becomes $u_i+cu_j$ and each dual column has $z_j$ replaced by
$z_j-cz_i$. Let $C(s)$ be the minimum number of paid additions to a successful state, allowing
free operations anywhere.

\begin{enthm}\label{e:main}
Suppose two labeled states have the same $u$ and $\A$, possibly with different parameter
dimensions. For any legal operation word and any
prescribed final order, the word passes the test from one state if and only if it passes from
the other. In particular their distances to the successful set agree. Quotienting the pairs
$(u,\A)$ further by free row permutations and sign changes still gives $C(s)$ and permits
recovery of an optimal word on the original input.
\end{enthm}
\noindent\emph{Proof.}\enspace
By~\eqref{e:action}, applying the same operation sends equal pairs to equal pairs. Induct on
word length, then use~\eqref{e:minor} on each prefix of the final order. This proves the
first assertion and preserves the paid length of every successful word. Free operations can
also change row labels and require a lifting argument.

\newpage
\section{Free operations, parameter changes, and recovery of a word}
Let $Q$ be a signed permutation matrix. Suppose the pair $(u,\A)$ of a second representative
is obtained from that of a first by $Q$. An addition $E=I+ce_i e_j^\top$ on the second has
the same effect on the pair as first applying $Q^{-1}EQ$ to the first and then applying $Q$.
The conjugate is again a single legal addition, with possibly different indices and coefficient.
Free operations stay in one class. A passing order on the second representative pulls back to
one on the first. These statements use only the action on pairs and apply to representatives
with different parameter presentations.

Give the quotient graph a directed unit edge for each paid addition, allowing loops. A physical
execution projects to a walk of the same paid length. Conversely, the conjugation rule lifts each edge from any
chosen physical representative, and a passing order can be chosen by a final free permutation.
The distance to the successful classes is therefore $C(s)$. In breadth-first search, store
one labeled state actually reached from the input for each new class, together with its
parent class and the addition used. Reversing the list of additions along the parent chain
recovers a shortest word on the input. This proves Theorem~\ref{e:main}.\qed

\medskip\noindent\textbf{Recognizing the same affine image.}
A common change of parameters $a\mapsto Ta+t$ acts by
\begin{equation}\label{e:frame}\frameeq\end{equation}
It preserves the affine image. Conversely, for fixed parameter dimension $d$, equality
$b+\im W=b'+\im W'$ gives $\im W=\im W'$ and $b'-b=Wt$ for some $t$. Map $\ker W'$
isomorphically onto $\ker W$, and identify complements of these kernels through their common
image. The resulting invertible map $T$ satisfies $W'=WT$. This includes dependent slope
columns and $d=0$.

Presentations of different dimensions can be compared directly. In homogeneous coordinates,
put
\begin{equation}\label{e:hom}\homoeq\end{equation}
The slice at last coordinate $1$ spans $H$. Differences of points in that slice give all of
$\im W\times\{0\}$, and the slice also contains $(b,1)$. Thus, with row labels fixed, a
complete key for $(u,\A)$ consists of $u$ and the nonzero rows of the reduced row-echelon
form of the transposed block matrix in~\eqref{e:hom}. For the free quotient, take the
lexicographically least key of $Qu$ and $\operatorname{diag}(Q^{-\top},1)H$ over all signed
permutations $Q$.

With one parameter, \eqref{e:frame} consists of the six maps $a\mapsto\pm a+t$. For each one,
replace every transformed row triple $(u_i,w_i,b_i)$ by the smaller of the triple and its
negative, and sort the list with multiplicities. The least of the six lists is the key used
in the census. In particular it retains the dual data in rows with $u_i=0$.

\medskip\noindent\textbf{Prescribed endpoints.}
Take $u=(1,0)^\top$, $W=(0,1)^\top$, and use $b=(1,0)^\top$ in one state and $b'=(1,1)^\top$
in the other. Both images are $\{(1,a)^\top:a\in K\}$, so the states give the same quotient
vertex. The support of $u$ fixes the first row, however, and sign changes cannot turn the
second constant entry from $0$ into $1$. There is no zero-cost path between these prescribed
labeled states.

\newpage
\section{Exact enumeration and scope}
The predecessor search uses a multiset of complete rows modulo signs. Its shortest-path
algorithm and recovery of a word were verified in Lean.~\cite{eorbits,elean} A separate
charge-only criterion decides the sublevels through two additions.~\cite{echarge} The present
quotient retains the primal column and the full affine image of the dual column, and its
distance theorem holds at every depth.

All three columns below include the free row operations. The other identifications are
parameter translations in the middle column and all invertible affine changes in the last.
For the five-row input, shell $k$ consists of classes whose shortest distance from the start
in the corresponding quotient graph is exactly $k$.
\counttable{Objects counted}{Free rows}{With translations}{With all affine maps}
  {All normalized three-row inputs}{Five rows, shell }

For three rows and one parameter there are $26\cdot9\cdot9=2\,106$ normalized inputs.
There are $26$ choices for nonzero $u$, and the separate dot-product conditions leave
$9$ choices each for $W$ and $b$. The verifier runs a multi-source $0$--$1$ breadth-first
search on labeled states. It includes explicit free swaps, sign changes and paid additions,
and evaluates the determinants of the actual prefix submatrices of $I-uz^\top$.
The quotient search tests $u_B^\top W_B=0$ and $u_B^\top b_B\ne1$ on each prefix $B$,
since a nonzero linear functional takes every value in $K$.
All three quotient distances agree with the labeled distance on each input. There are
$882$ inputs of cost zero, $1\,080$ of cost one, and $144$ of cost two.

The five-row input from the predecessor has $u=(1,1,1,1,1)^\top$,
$W=(1,1,1,1,-1)^\top$, and $b=(1,1,1,1,0)^\top$. The first successful shell is shell
three. With zero-based labels the search recovers the word
\[
  \fivewitness
\]
After execution $(u_i z_i)_i=(0,0,0,0,1)$ for every parameter. All proper leading minors in
the original order are therefore $1$. The third shell has 86 fewer classes. Runtime also
depends on the cost of computing a key.

For two rows, the verifier separately checks all $819$ presentations $(W,b)$ with
$d=0,1,2$. The key in~\eqref{e:hom} and explicit enumeration of affine points give the
same $22$ classes. This image-key check uses unrestricted $(W,b)$. All computations use
the Python 3.12.10 standard library and can be reproduced from the package directory with
\texttt{python -B verify.py}.

\medskip\noindent\textbf{Formal verification and an open question.}
Lean 4.24.0 verifies path lifting at every length, optimal witness recovery in Theorem~\ref{e:main},
and the five-row optimum of three. It also verifies the five-row closed balls of radii two and three.
The row quotient grows from $47$ to $213$ classes and the affine quotient from $42$ to $122$.
Subtraction gives the shell counts $166$ and $80$. Python reproduces the full table.
The relative runtime of the encodings for several parameters remains open.
\sources{e}
\end{document}
